We prove that every non-trivial zero of the Riemann zeta function, and of every Dirichlet $L$-function,
has real part $1/2$. The argument transfers Stepanov's method for curves over finite fields to the
integers.

A Landau-type theorem on windows reduces the problem to one-sided bounds, with square-root error,
for Riesz-weighted prime sums. An exact orthogonality identity on auxiliary abelian ``host'' fields
turns these into upper bounds for weighted counts of primes that are inert in an imaginary quadratic
field and lie in prescribed Frobenius classes. These counts are bounded by comparing two quantities
attached to lattices built from chains of theta lifts between totally definite unitary groups:
their divisibility at inert primes, and archimedean slope and covolume bounds. The chains have length
comparable to the square root of the length of the prime sum, and they play the role of the auxiliary
multiplicity in Stepanov's method. Galois-cohomological obstructions are controlled by bounds for
Tate--Shafarevich groups of cyclotomic twists over the host fields.

The complete proof has been formalized in Lean~4 with Mathlib, and the formal proof has been checked
independently by several proof kernels. Part~I of this paper explains the architecture of the argument.
It gives complete proofs of the analytic reductions, and it states precisely each arithmetic input,
together with an outline of its proof. Part~II gives the proofs of the arithmetic inputs, converted from
the formal proof into mathematical language; its deepest lemmas are cited by the names of their formal
counterparts.
